\section{投稿者の意図に基づく定義の影響}
\label{sec_intent}
本章では，第4版と第5版のお題文の定義の下で行った評価から，意図に基づく定義が教師，生徒，全データのラベルに与えた影響を示す．

\subsection{評価の方法}
お題文ありの画像は全データの約8\%と少ないため，無作為に抽出した目視ラベルだけでは正例が不足する．
そこで，パッケージの予測で代表の画像を層に分け，層ごとに無作為に抽出して目視ラベルを付け，層の重みで全データの値を推定した．
層 $ h $ に属する代表の画像の数を $ N_h $，そのうち目視ラベルを付けた画像の数を $ n_h $ とすると，層 $ h $ の画像 $ i $ の重みは $ w_i = N_h / n_h $ である．
目視ラベルを $ y_i \in \{0, 1\} $，予測を $ \hat{y}_i \in \{0, 1\} $ とすると，全データの適合率 $ P $ と再現率 $ R $ の推定値は(\ref{eq_weighted_pr})で表される．

\begin{equation}
    P = \frac{\sum_{i} w_i y_i \hat{y}_i}{\sum_{i} w_i \hat{y}_i}, \quad
    R = \frac{\sum_{i} w_i y_i \hat{y}_i}{\sum_{i} w_i y_i}
    \label{eq_weighted_pr}
\end{equation}

\noindent
F1 は $ P $ と $ R $ の調和平均であり，95\%区間は層ごとの復元抽出によるブートストラップで求めた．
評価には，第4版の目視ラベルを付けた1,020枚と，第5版のお題文ありの評価のために新たに抽出した600枚を用いた．
1,020枚は，評価用データの640枚と，実データの無作為抽出160枚および層別抽出220枚からなる．
600枚は，パッケージの形式の予測とお題文ありの生徒の確率で分けた9層からの抽出であり，お題文ありの正例を多く含む．

\subsection{教師への確認の質問}
第4版の段階で，教師は元の素材の文字をお題文とする誤りが多かった．
定義の説明に「元からある文字は大きくてもお題文ではない」と書き足すと，実データでお題文と誤る画像はかえって13枚から18枚に増えた．
一方，教師がお題文とした画像に限って「投稿者が回答者に向けて書いた問い，指示，題目の句があるか」を別の質問で確かめると，実データでの誤りは13枚から4枚に減り，実データでの教師のお題文の F1 は0.769から0.873に上がった．

しかし，この確認の質問は，お題文と通常の区別には働いていなかった．
確認の質問を目視ラベルのある1,020枚の全てに問うと，お題文の適合率は0.20，再現率は1.00であった．
「はい」と答えた割合は，目視でお題文の画像で96/96，穴埋めで192/259，カルタで50/69，通常で141/596であった．
確認の質問は，題であるかよりも，投稿者が書き加えたらしい文の有無に反応していた．
また，生徒がお題文とした画像に限ると，確認の質問は目視ラベルのある画像で一度も「いいえ」と答えず，全データでも対象の567,345枚のうち0.7\%の判定しか変えなかった．

\subsection{第5版のお題文ありと VLM}
第5版でお題文ありを独立の属性としたため，穴埋めとカルタの画像328枚にお題文ありの目視ラベルを付け直したところ，30枚がお題文ありであった．
表\ref{tab_qwen_v5}に，この正解に対する教師の F1 を示す．
評価用は1,020枚のうちの最終評価用の320枚，実データは無作為抽出と層別抽出の380枚の重み付きの値である．

\begin{table}[t]
    \centering
    \caption{第5版のお題文ありに対する教師の F1 と推論の速さ}
    \label{tab_qwen_v5}
    \small
    \begin{tabular}{lccc}
        \toprule
        モデルと指示 & 枚/秒 & 評価用 & 実データ \\
        \midrule
        8B，確認の質問 & 18.9 & 0.43 & 0.39 \\
        8B，定義と手順 & 12.7 & 0.57 & 0.56 \\
        32B，確認の質問 & 9.7 & 0.52 & 0.53 \\
        32B，定義と手順 & 6.0 & 0.76 & 0.75 \\
        \bottomrule
    \end{tabular}
\end{table}

8B のモデルは，定義と判断の手順を詳しく与えてもお題文ありと答えすぎた．
誤りの多くは，空欄を含む文そのもの，テレビの字幕，漫画の台詞をお題文とするものであった．
32B のモデルに定義と手順を与えると，除外の条件をおおむね守れたが，F1 は0.75前後にとどまった．
32B のモデルを全データに適用すると約12.8日を要するため，その判定を4.5万枚の教師ラベルとして生徒に学習させ，パッケージ v5.3.0 に組み込んだ．

\subsection{第5版の評価と残った誤り}
600枚のうち層別抽出の420枚で，パッケージ v5.3.0 のお題文ありの F1 は0.831であり，お題文ありの全データの割合は7.9\%と推定された．
同じ判断者による90枚の付け直しの一致は，一致率0.978，Cohen の $ \kappa $\cite{ref_kappa} 0.95 であった．
不一致の2枚は，話者を付けた台本と書き写した歌詞であり，いずれも付け直しの時点で境界の規則を定めたことによる差であった．

残った誤りは，文字の出所の判断に集中した．
穴埋めやカルタと予測され，お題文ありの生徒の確率が0.3以上の層では，目視でお題文ありの画像は35\%にとどまった．
この層のお題文なしの画像の多くは，空欄を含む1文だけの画像や，空欄を含む歌詞，記事，名札の見出し付きの問題文であった．
また，形式が通常と予測された画像には，全データで約3.0万枚のお題文ありの見逃しがあると推定された．
生徒の確率が0.7以上の通常の画像でも，お題文ありは56\%であった．
これらの画像の多くは，写真の上下の帯や写真に重ねた文字の画像であり，人物の台詞の字幕と，投稿者が加えた問いや状況の文が同じ見た目になる．
32B のモデルも，この中のお題文なしの74枚のうち33枚をお題文ありと答えた．

以上から，意図に基づく定義の下では，教師と生徒の誤りの多くが，画像から観測できない文字の出所の判断に集中していたといえる．
境界の規則を書き足すことで目視ラベルの再現性は保たれたが，同じ見た目の画像に異なるラベルを求めるため，教師と生徒が学ぶべき区別そのものが画像に現れない．
